\textbf{Lösung zu Klausurvorbereitung 3}

\begin{loesung}
\textbf{a)} Euklidischer Algorithmus
\begin{align*}
a &= q_0\, b + r_1\\
b &= q_1\, r_1 + r_2\\
&\cdots\\
r_{n-1} &= q_n\, r_n \qquad \Rightarrow \quad r_n = \ggT(a, b)
\end{align*}
\begin{align*}
2406 &= 3 \cdot 654 + 444\\
654 &= 1 \cdot 444 + 210\\
444 &= 2 \cdot 210 + 24\\
210 &= 8 \cdot 24 + 18\\
24 &= 1 \cdot 18 + 6\\
18 &= 3 \cdot 6 \quad \Rightarrow \quad r_n = \ggT(a, b) = 6
\end{align*}

\textbf{b)} Bestimmen Sie eine ganzzahlige Lösung der Gleichung
\[ 2406\,x - 654\,y = 24 \]
24 ist ein Vielfaches von $\ggT(2406, 654)$, diese Gleichung ist daher lösbar.

\begin{enumerate}
\item Lösung: EA rückwärts
\begin{align*}
24 &= 444 - 2 \cdot 210\\
24 &= 444 - 2 \cdot (654 - 444)\\
24 &= 3 \cdot 444 - 2 \cdot 654\\
24 &= 3 \cdot (2406 - 3 \cdot 654) - 2 \cdot 654\\
24 &= 3 \cdot 2406 - 11 \cdot 654
\end{align*}
\item Lösung: Mit der Matrix $Q$
\begin{align*}
Q &= \begin{pmatrix} 3 & 1\\ 1 & 0\end{pmatrix}
\begin{pmatrix} 1 & 1\\ 1 & 0\end{pmatrix}
\begin{pmatrix} 2 & 1\\ 1 & 0\end{pmatrix}
\begin{pmatrix} 8 & 1\\ 1 & 0\end{pmatrix}
\begin{pmatrix} 1 & 1\\ 1 & 0\end{pmatrix}
\begin{pmatrix} 3 & 1\\ 1 & 0\end{pmatrix}\\
&= \begin{pmatrix} 4 & 3\\ 1 & 1\end{pmatrix}
\begin{pmatrix} 17 & 2\\ 8 & 1\end{pmatrix}
\begin{pmatrix} 4 & 1\\ 3 & 1\end{pmatrix}
= \begin{pmatrix} 4 & 3\\ 1 & 1\end{pmatrix}
\begin{pmatrix} 74 & 19\\ 35 & 9\end{pmatrix}
= \begin{pmatrix} 401 & 103\\ 109 & 28\end{pmatrix}
\end{align*}
\begin{align*}
% KORRIGIERT: Begründung für den Übergang zur nächsten Zeile ergänzt.
\operatorname{Det} Q = 401 \cdot 28 - 109 \cdot 103 &= 1 && | \cdot 6 \quad \text{(es ist } 6 \cdot 401 = 2406,\ 6 \cdot 109 = 654\text{)}\\
2406 \cdot 28 - 654 \cdot 103 &= 6 && | \cdot 4\\
2406 \cdot 112 - 654 \cdot 412 &= 24
\end{align*}
\end{enumerate}

% KORRIGIERT: Original: "Für alle weiteren Lösungen dieser Gleichung muss gelten 2406x - 654y = 0" (gilt nur für die Differenz zweier Lösungen).
\textbf{c)} Ist $(x, y)$ eine weitere Lösung, so muss für die Differenz zur Lösung $(3, 11)$ gelten
\[ 2406 \cdot (x - 3) - 654 \cdot (y - 11) = 0 \]
Es ist
\[ 2406 \cdot \frac{654}{6} - 654 \cdot \frac{2406}{6} = 0 \]
also
\[ 2406 \cdot 109 - 654 \cdot 401 = 0 \]
% KORRIGIERT: Begründung der Vollständigkeit und "t \in \Z" ergänzt.
Da $\ggT(401, 109) = 1$ ist, sind dies bis auf ganzzahlige Vielfache alle Lösungen von $2406\,x - 654\,y = 0$.
Damit erhält man für die allgemeine Lösung von $2406\,x - 654\,y = 24$
\[ x = 3 + 109\,t, \qquad y = 11 + 401\,t, \qquad t \in \Z \]
\end{loesung}
